% Bonding parameter reference -- all 31 parameters, grouped by the six
% Parameter Edit screens.
%
% DERIVED FROM FIRMWARE. Refresh against:
%   Firmware/Core/Src/configuration_parameter_catalog.cpp  (screen, mode mask, scale)
%   Firmware/Core/Inc/configuration.h                      (min/max/step, defaults)
%   Firmware/Core/Src/menu.cpp                             (display names)
%   Firmware/Core/Inc/bonder_config.hpp                    (field -> default mapping)
%
% Min/max/step/default are DISPLAY values (raw * scale). Power, energy and
% timing are stored in SI (W, J, s) and shown in mW, mJ, ms.
%
% Mode column: S = Semi Auto, M = Manual, T = Table Tear, L = Lange Cplr.
% A parameter not available in the active mode is hidden from its screen.

\newcommand{\ptable}[1]{%
  \begin{center}\small
  \begin{tabular}{@{}L{28mm}L{10mm}R{16mm}R{16mm}R{13mm}R{16mm}L{16mm}@{}}
  \toprule
  \textbf{Name} & \textbf{Unit} & \textbf{Minimum} & \textbf{Maximum} &
  \textbf{Step} & \textbf{Default} & \textbf{Modes} \\
  \midrule
  #1
  \bottomrule
  \end{tabular}
  \end{center}}

% ------------------------------------------------------------- screen 1/6 --

\subsection*{Screen 1/6 --- First bond}

\ptable{
\menuitem{Search 1}  & mm & 0.00 & 20.00 & 0.01 & 3.50 & S \ T \ L \\
\menuitem{Power 1}   & mW & 0    & 1000  & 10   & 500  & all \\
\menuitem{Energy 1}  & mJ & 0    & 1000  & 1    & 50   & all \\
\menuitem{Force 1 G} & g  & 0    & 150   & 1    & 35   & all \\
}

\begin{description}[leftmargin=!,labelwidth=26mm,style=nextline,font=\normalfont\ttfamily]
\item[Search 1] Height the tool descends to under tracking force before it
  begins its controlled search for the pad. Lower it to shorten cycle time;
  raise it if the tool risks striking the workpiece on the way in.
\item[Power 1] Electrical drive power for the first weld. Sets ultrasonic
  \emph{amplitude} --- how hard the tool scrubs.
\item[Energy 1] Total ultrasonic energy delivered into the first weld. Sets the
  \emph{dose}. Together with \menuitem{Power 1} this implies the weld time,
  which is capped by \menuitem{Bond Timeout}.
\item[Force 1 G] Bond force held during the first weld.
\end{description}

% ------------------------------------------------------------- screen 2/6 --

\subsection*{Screen 2/6 --- Loop formation}

\ptable{
\menuitem{Stepback}    & mm & 0.00  & 5.00  & 0.01 & 5.00 & S \ T \ L \\
\menuitem{Kink Height} & mm & -2.00 & 5.00  & 0.01 & 4.00 & all \\
\menuitem{Reverse}     & mm & 0.00  & 5.00  & 0.01 & 0.00 & all \\
\menuitem{Loop Height} & mm & 0.00  & 20.00 & 0.01 & 5.00 & all \\
}

\begin{description}[leftmargin=!,labelwidth=26mm,style=nextline,font=\normalfont\ttfamily]
\item[Stepback] Y position the stage moves to while preparing the second bond.
\item[Kink Height] Intermediate Z height the head rises to immediately after
  the first weld, forming the kink at the heel of the bond. Reached before the
  tail is fed.
\item[Reverse] Distance the Y stage backs off, measured from wherever it
  currently sits. Zero by default. Semi Auto, Manual and Table Tear apply it
  after the kink, before the loop rises; Lange Coupling instead applies it
  during the second-bond descent, producing the drag that mode is used for.
\item[Loop Height] Z apex between the two bonds --- the top of the wire loop.
\end{description}

% ------------------------------------------------------------- screen 3/6 --

\subsection*{Screen 3/6 --- Second bond}

\ptable{
\menuitem{Search 2}  & mm & 0.00 & 20.00 & 0.01 & 3.50 & S \ T \ L \\
\menuitem{Power 2}   & mW & 0    & 1000  & 10   & 700  & all \\
\menuitem{Energy 2}  & mJ & 0    & 1000  & 1    & 70   & all \\
\menuitem{Force 2 G} & g  & 0    & 150   & 1    & 35   & all \\
\menuitem{Tail}      & mm & 0.00 & 10.00 & 0.01 & 4.50 & S \ M \ L \\
\menuitem{Tear}      & mm & 0.00 & 10.00 & 0.01 & 4.50 & S \ M \ L \\
}

\begin{description}[leftmargin=!,labelwidth=26mm,style=nextline,font=\normalfont\ttfamily]
\item[Search 2] As \menuitem{Search 1}, for the second bond.
\item[Power 2] Ultrasonic amplitude for the second weld. Defaults higher than
  \menuitem{Power 1} because the second bond is made through the wire.
\item[Energy 2] Ultrasonic dose for the second weld.
\item[Force 2 G] Bond force held during the second weld.
\item[Tail] T-axis travel that feeds the wire tail after the first bond. This
  is a \emph{relative} move added to wherever the T axis currently sits, not an
  absolute coordinate. In Table Tear mode the equivalent is
  \menuitem{Table Tail}.
\item[Tear] T-axis travel that tears the wire after the second weld, again
  relative to the axis's current position. Because it follows \menuitem{Tail},
  the endpoint depends on where the tail move left the axis.
\end{description}

% ------------------------------------------------------------- screen 4/6 --

\subsection*{Screen 4/6 --- Heights and positions}

\ptable{
\menuitem{Reset Height} & mm & 0.00  & 20.00 & 0.01 & 8.00  & all \\
\menuitem{Overtravel}   & mm & -2.00 & 0.00  & 0.01 & -0.30 & S \ T \ L \\
\menuitem{Second Z Hgt} & mm & 0.00  & 20.00 & 0.01 & 2.50  & T \\
\menuitem{Table Tail}   & mm & 0.00  & 5.00  & 0.01 & 2.00  & T \\
\menuitem{Table Tear}   & mm & 0.00  & 5.00  & 0.01 & 3.50  & T \\
}

\begin{description}[leftmargin=!,labelwidth=26mm,style=nextline,font=\normalfont\ttfamily]
\item[Reset Height] Parked Z height. The machine returns here at the end of
  every cycle and is driven here before every cycle begins.
\item[Overtravel] Z target \emph{below} the pad surface, so the descent is a
  direction to push in rather than a position to arrive at --- the workpiece
  stops the tool. Always zero or negative. Make it less negative if the tool
  digs in; more negative if it fails to reach low pads.
\item[Second Z Hgt] Table Tear mode only. Z height held while the Y stage forms
  and tears the tail.
\item[Table Tail] Table Tear mode only. Y position that forms the tail.
\item[Table Tear] Table Tear mode only. Y position that tears the wire.
\end{description}

% ------------------------------------------------------------- screen 5/6 --

\subsection*{Screen 5/6 --- Timing}

\ptable{
\menuitem{Bond Timeout} & ms & 0 & 60000 & 100 & 5000 & all \\
\menuitem{Contact Stl}  & ms & 0 & 5000  & 10  & 100  & all \\
\menuitem{Cooling}      & ms & 0 & 5000  & 10  & 10   & all \\
\menuitem{Tail Delay}   & ms & 0 & 5000  & 10  & 1000 & S \ M \ L \\
\menuitem{Tear Stabil}  & ms & 0 & 5000  & 10  & 100  & T \\
}

\begin{description}[leftmargin=!,labelwidth=26mm,style=nextline,font=\normalfont\ttfamily]
\item[Bond Timeout] Safety ceiling on ultrasonic delivery. If the requested
  energy has not been delivered within this time the weld is abandoned and the
  machine reports \msg{LOW BONDING POWER}. This is a fault limit, not a
  process setting --- it should never be reached in normal bonding.
\item[Contact Stl] Dwell after touchdown and again after the force ramp,
  letting the force settle before ultrasonics start.
\item[Cooling] Dwell after the weld, before the force is released, letting the
  bond solidify.
\item[Tail Delay] Dwell after the tear, before the tail is restored.
\item[Tear Stabil] Table Tear mode only. Dwell after the tear, before the axes
  are restored.
\end{description}

% ------------------------------------------------------------- screen 6/6 --

\subsection*{Screen 6/6 --- Force and ultrasonic setup}

\ptable{
\menuitem{Constant G}  & g   & 0    & 150  & 1   & 15   & all \\
\menuitem{Tracking G}  & g   & 0    & 150  & 1   & 30   & all \\
\menuitem{Scan Start}  & kHz & 50.0 & 70.0 & 0.1 & 58.0 & all \\
\menuitem{Scan Stop}   & kHz & 50.0 & 70.0 & 0.1 & 62.0 & all \\
\menuitem{Scan Points} & ---  & 8    & 32   & 1   & 16   & all \\
\menuitem{Tail Power}  & mW  & 0    & 1000 & 10  & 200  & S \ M \ L \\
\menuitem{Tail Energy} & mJ  & 0    & 1000 & 1   & 200  & S \ M \ L \\
}

\begin{description}[leftmargin=!,labelwidth=26mm,style=nextline,font=\normalfont\ttfamily]
\item[Constant G] Light holding force carried during descent and between
  process steps --- enough to keep the lever controlled, not enough to bond.
\item[Tracking G] Force loaded while the head searches for the pad.
\item[Scan Start] Low edge of the impedance sweep used to find the
  transducer's resonance.
\item[Scan Stop] High edge of that sweep. \menuitem{Scan Start} and
  \menuitem{Scan Stop} must bracket the actual resonance; confirm with the
  \keycap{TEST} button.
\item[Scan Points] Number of frequencies measured across the sweep. More points
  give a better fit at the cost of a slightly longer scan.
\item[Tail Power] Ultrasonic amplitude for the tail assist --- a brief burst
  applied while the tail is re-formed after the tear. Also the power used by
  the \keycap{TEST} function.
\item[Tail Energy] Ultrasonic dose for the tail assist, and for \keycap{TEST}.
\end{description}

\derivedfrom{\texttt{configuration\_parameter\_catalog.cpp},
             \texttt{configuration.h}, \texttt{menu.cpp}}
